% !TeX root = ../presentacion.tex
% ===========================================================================
%  Láminas de respaldo
%
%  Van después de \appendix: no cuentan en la numeración ni en la barra de
%  progreso. Son las respuestas a las preguntas que uno espera, ya
%  maquetadas. Se llega a ellas saltando por número desde el visor.
%
%  Anota aquí el número de cada lámina cuando la presentación esté cerrada,
%  y llévalo en un papel: buscar una lámina de respaldo pasando páginas
%  delante del público arruina el efecto.
% ===========================================================================

\begin{frame}{Índice de respaldo}
  \guion{Para el ponente. Apunta el número de lámina de cada fila cuando la
  presentación esté congelada.}

  \begin{itemize}\footnotesize
    \item Vistas C4 restantes: componentes de cada servicio y microfrontend
    \item Despliegue (\vista{04}) y el papel del \id{edge}
    \item Registro completo de decisiones (ADR)
    \item Diccionario de datos y DDL
    \item Traducción de reglas de negocio a códigos HTTP
    \item Estrategia de pruebas
    \item Por qué DDD y dónde están los agregados
  \end{itemize}
\end{frame}

\laminafiguratitulo{Despliegue (\vista{04})}{despliegue}%
  {Dos redes Docker. Solo el \id{edge} toca las dos.}

\laminafiguratitulo{Componentes del \id{shell}}{componentes-shell}%
  {Cómo el host resuelve los remotes y comparte la sesión.}

\laminafiguratitulo{Componentes de \id{ms-canchas}}{componentes-canchas}{}

\laminafiguratitulo{Componentes de \id{ms-usuarios}}{componentes-usuarios}{}

\laminafiguratitulo{Componentes de \id{ms-reportes}}{componentes-reportes}%
  {El único servicio sin base de datos: compone por HTTP.}

\laminafiguratitulo{Componentes del \id{api-gateway}}{componentes-gateway}%
  {Dónde se verifica la credencial antes de rutear.}

\begin{frame}{Reglas de negocio $\rightarrow$ respuestas HTTP}
  \guion{Respuesta preparada para «¿y qué le devuelve al cliente cuando la
  regla no se cumple?». Rellenar con lo que hace el servicio de verdad.}

  \begin{center}
    \small
    \begin{tabular}{@{}llp{5.2cm}@{}}
      \toprule
      \textbf{Situación} & \textbf{Código} & \textbf{Regla}\\
      \midrule
      Bloque ya ocupado          & \id{409} & \rn{02}\\
      Cancelar una reserva ajena & \id{403} & \rn{03}\\
      Cancelar una reserva pasada& \id{422} & \rn{04}\\
      Tope de reservas alcanzado & \id{422} & \rn{06}\\
      \ph{Completar con el resto}& \ph{---} & \\
      \bottomrule
    \end{tabular}
  \end{center}
  {\footnotesize\color{textoSuave} Verifica cada código contra el
  controlador antes de exponer: un número inventado aquí se nota enseguida.}
\end{frame}

\begin{frame}{Estrategia de pruebas}
  \ph{Pirámide de pruebas del proyecto: unitarias de dominio, de integración
  con base real, y de contrato. Cifras concretas si las tenéis.}
\end{frame}

\begin{frame}{Preguntas que esperamos}
  \guion{Ensayad las respuestas en voz alta. Una respuesta pensada en el
  momento suena a que no se pensó nunca.}

  \begin{itemize}\footnotesize
    \item ¿Por qué microservicios para un sistema de este tamaño?
    \item ¿Qué pasa si \id{ms-canchas} está caído cuando alguien reserva?
    \item ¿Por qué no una cola de mensajes entre servicios?
    \item ¿Cómo escalaríais si mañana hubiera cien clubes?
    \item ¿Module Federation no acopla los remotes al shell por la versión de React?
    \item ¿Cómo evitáis que un remote roto tumbe toda la aplicación?
    \item \ph{La que os dé más miedo}
  \end{itemize}
\end{frame}
